\documentclass[aspectratio=169,11pt]{beamer}
\usetheme{default}
\setbeamertemplate{navigation symbols}{}
\setbeamertemplate{footline}{%
  \hfill\insertframenumber/\inserttotalframenumber\hspace*{4pt}\vspace{2pt}%
}
\usepackage{lmodern}
\usepackage[most]{tcolorbox}
\usepackage{listings}
\usepackage{tikz}
\usetikzlibrary{arrows.meta,positioning,shapes.geometric,fit,backgrounds}
\usepackage{booktabs}
\usepackage{hyperref}

\definecolor{lcgreen}{HTML}{2E7D32}
\definecolor{lcblue}{HTML}{1565C0}
\definecolor{lcorange}{HTML}{EF6C00}
\definecolor{lcred}{HTML}{C62828}
\definecolor{codebg}{HTML}{F5F5F5}
\definecolor{docgray}{gray}{0.55}

\newtcolorbox{conceptbox}[1][]{colback=yellow!20,colframe=yellow!60!black,title=#1,fonttitle=\bfseries,top=2pt,bottom=2pt,left=4pt,right=4pt}
\newtcolorbox{warningbox}[1][]{colback=red!15,colframe=red!60!black,title=#1,fonttitle=\bfseries,top=2pt,bottom=2pt,left=4pt,right=4pt}
\newtcolorbox{tipbox}[1][]{colback=green!10,colframe=green!50!black,title=#1,fonttitle=\bfseries,top=2pt,bottom=2pt,left=4pt,right=4pt}
\newtcolorbox{codebox}[1][]{colback=blue!10,colframe=blue!40!black,title=#1,fonttitle=\bfseries,top=2pt,bottom=2pt,left=4pt,right=4pt}

\lstset{
  basicstyle=\ttfamily\scriptsize,
  backgroundcolor=\color{codebg},
  breaklines=true,
  frame=single,
  rulecolor=\color{blue!40!black},
  keywordstyle=\color{lcblue}\bfseries,
  stringstyle=\color{lcgreen},
  commentstyle=\color{docgray}\itshape,
  language=Python,
  showstringspaces=false,
  tabsize=2,
  xleftmargin=2pt,
  xrightmargin=2pt,
  aboveskip=2pt,
  belowskip=2pt,
}

\newcommand{\docref}[1]{{\scriptsize\color{docgray}[Docs: #1]}}

\title{LangGraph v1.0.x}
\subtitle{Section 01: State, Nodes \& Edges}
\author{Agentic AI Foundations}
\date{March 2026}

\begin{document}

\begin{frame}
\titlepage
\end{frame}

\section{Mental Model}
\begin{frame}[t]{Three Building Blocks of Every LangGraph App}
\setlength{\itemsep}{1pt}
\vspace{0.1cm}

\begin{conceptbox}[State + Nodes + Edges = Graph]
\setlength{\itemsep}{1pt}
\begin{itemize}
  \item \textbf{State:} The shared data structure (TypedDict)
  \item \textbf{Nodes:} Functions that compute and update state
  \item \textbf{Edges:} Routes that connect nodes together
\end{itemize}
\end{conceptbox}

\vspace{0.15cm}

Every LangGraph application starts here. Master these three, and you can build any workflow.

\vspace{0.15cm}

\begin{tikzpicture}[node distance=1.2cm, scale=0.9, every node/.style={font=\scriptsize}]
  \node[draw, fill=lcblue!30, rounded rectangle, minimum width=1.5cm] (s) {State};
  \node[draw, fill=lcgreen!30, rounded rectangle, right=of s, minimum width=1.5cm] (n) {Node};
  \node[draw, fill=lcorange!30, rounded rectangle, right=of n, minimum width=1.5cm] (e) {Edge};
  \path[->] (s) -- (n);
  \path[->] (n) -- (e);
\end{tikzpicture}
\end{frame}

\section{State Deep Dive}
\begin{frame}[fragile,t]{State: TypedDict Schema}
\vspace{0.05cm}

\begin{conceptbox}[State is a TypedDict]
State defines the \textit{shape} of data flowing through your graph. It's just a Python dict with types.
\end{conceptbox}

\vspace{0.1cm}

\begin{codebox}[Example State]
\begin{lstlisting}
from typing_extensions import TypedDict

class State(TypedDict):
    user_input: str
    messages: list
    decision: str
\end{lstlisting}
\end{codebox}

\begin{tipbox}[Why]
Type hints enable validation and IDE autocomplete.
\end{tipbox}
\end{frame}

\begin{frame}[fragile,t,shrink=2]{State Fields: Defining Your Schema}
\vspace{0.05cm}

{\footnotesize
\begin{tabular}{p{2.8cm}p{6cm}}
\toprule
\textbf{Pattern} & \textbf{Use Case} \\
\midrule
\texttt{user\_input: str} & Single string field \\
\texttt{messages: list} & Chat history, append-only \\
\texttt{score: float} & Numeric decision \\
\texttt{context: dict} & Nested data \\
\texttt{Optional[str]} & Field can be None \\
\bottomrule
\end{tabular}
}

\vspace{0.08cm}

\begin{warningbox}[No Inheritance]
\texttt{TypedDict} does NOT support inheritance. Use composition.
\end{warningbox}

\vspace{0.05cm}

{\footnotesize State fields are the interface. Good design = good graph.}
\end{frame}

\section{Reducers}
\begin{frame}[fragile,t]{Reducers: How State Updates Work}
\setlength{\itemsep}{1pt}
\vspace{0.1cm}

\begin{conceptbox}[Default Behavior: Overwrite]
When a node returns \texttt{\{"field": new\_value\}}, the state field is \textit{replaced}.
\end{conceptbox}

\vspace{0.15cm}

\begin{codebox}[Default Merge: Overwrite]
\begin{lstlisting}
# Node returns partial state
return {"score": 0.85}  # replaces state["score"]
\end{lstlisting}
\end{codebox}

\vspace{0.15cm}

\begin{tipbox}[Why Reducers Matter]
For \textbf{lists} (messages, history), you usually want to \textit{append}, not replace. That's where reducers come in.
\end{tipbox}
\end{frame}

\begin{frame}[fragile,t,shrink=5]{Annotated: Specifying Reducers}
\vspace{0.1cm}

\begin{conceptbox}[Annotated Pattern]
Use \texttt{Annotated[type, reducer\_function]} to specify custom merge behavior.
\end{conceptbox}

\vspace{0.1cm}

\begin{codebox}[Example: List Append]
\begin{lstlisting}
from typing_extensions import TypedDict, Annotated
import operator

class State(TypedDict):
    messages: Annotated[list, operator.add]
    score: float
\end{lstlisting}
\end{codebox}

\vspace{0.1cm}

\begin{tipbox}[What happens]
\texttt{operator.add} means: \texttt{state["messages"] = state["messages"] + new\_messages}. Appends!
\end{tipbox}

\vspace{0.1cm}

{\footnotesize When a node returns \texttt{\{"messages": [new\_msg]\}}, it appends, not replaces.}
\end{frame}

\section{Built-in Reducers}
\begin{frame}[fragile,t,shrink=2]{Built-in Reducers: operator.add}
\vspace{0.05cm}

\begin{tabular}{p{3.5cm}p{6.5cm}}
\toprule
\textbf{Reducer} & \textbf{Behavior} \\
\midrule
\texttt{operator.add} & \texttt{field + new\_value} (list concat) \\
\texttt{add\_messages} & Smart message appending (dedup, etc) \\
\texttt{Custom} & Define your own merging logic \\
\bottomrule
\end{tabular}

\vspace{0.15cm}

\begin{codebox}[operator.add Example]
\begin{lstlisting}
class State(TypedDict):
    items: Annotated[list, operator.add]

# Node returns {"items": [4, 5]}
# Result: state["items"] = [1, 2, 3] + [4, 5]
#       = [1, 2, 3, 4, 5]
\end{lstlisting}
\end{codebox}
\end{frame}

\begin{frame}[fragile,t,shrink=5]{add\_messages: The Chat Reducer}
\vspace{0.1cm}

\begin{conceptbox}[add\_messages: Most Important Reducer]
Built-in reducer from LangChain. Intelligently merges chat messages: handles duplicates, preserves order.
\end{conceptbox}

\vspace{0.1cm}

\begin{codebox}[Chat State Example]
\begin{lstlisting}
from langchain_core.messages import BaseMessage
from langgraph.graph import add_messages

class State(TypedDict):
    messages: Annotated[list[BaseMessage],
                        add_messages]
\end{lstlisting}
\end{codebox}

\vspace{0.1cm}

\begin{tipbox}[Why add\_messages?]
Handles HumanMessage, AIMessage, ToolMessage intelligently. Deduplicates, respects IDs.
\end{tipbox}

{\footnotesize Default for LLM workflows.}
\end{frame}

\section{Complex State}
\begin{frame}[fragile,t,shrink=3]{Complex State: Nested, Optional, Multiple}
\vspace{0.02cm}

\begin{codebox}[Advanced State Pattern]
\begin{lstlisting}
from typing_extensions import TypedDict, Annotated
from typing import Optional

class State(TypedDict):
    messages: Annotated[list, add_messages]
    user: Optional[dict]  # Optional field
    metadata: dict  # Nested dict
    scores: Annotated[list[float], operator.add]
\end{lstlisting}
\end{codebox}

\vspace{0.1cm}

\begin{tipbox}[Guidelines]
\setlength{\itemsep}{1pt}
\begin{itemize}
  \item Keep state \textit{flat} (avoid deep nesting)
  \item Use Annotated for lists/accumulators
  \item Optional for fields that may not exist yet
\end{itemize}
\end{tipbox}
\end{frame}

\section{Nodes}
\begin{frame}[fragile,t]{Nodes: Functions That Update State}
\vspace{0.1cm}

\begin{conceptbox}[Node Signature]
A node is a function: \texttt{(state: State) -> dict}. Takes state, returns partial updates.
\end{conceptbox}

\vspace{0.1cm}

\begin{codebox}[Simple Node]
\begin{lstlisting}
def process_input(state: State) -> dict:
    # Do computation
    result = state["user_input"].upper()
    # Return partial state update
    return {"decision": result}
\end{lstlisting}
\end{codebox}

\vspace{0.1cm}

\begin{warningbox}[Return Partial, Not Full]
Return ONLY the fields you update. The graph merges with existing state automatically.
\end{warningbox}
\end{frame}

\begin{frame}[fragile,t]{Node Return Values}
\vspace{0.1cm}

\begin{tabular}{p{4cm}p{7cm}}
\toprule
\textbf{Pattern} & \textbf{Effect} \\
\midrule
\texttt{\{field: value\}} & Update \textit{that} field (merge) \\
\texttt{\{field1: v1, field2: v2\}} & Update multiple fields \\
\texttt{\{\}} & Return empty dict (no update) \\
\bottomrule
\end{tabular}

\vspace{0.15cm}

\begin{codebox}[Examples]
\begin{lstlisting}
return {"score": 0.9}  # Update score only
return {"a": 1, "b": 2}  # Update a and b
return {}  # No state change
\end{lstlisting}
\end{codebox}

\vspace{0.15cm}

{\footnotesize The graph \textit{merges} your return dict into state. No need to copy full state.}
\end{frame}

\section{Node Patterns}
\begin{frame}[fragile,t]{Node Patterns: LLM Calls}
\vspace{0.05cm}

\begin{codebox}[LLM Node]
\begin{lstlisting}
from langchain_openai import ChatOpenAI
from langchain_core.messages import HumanMessage

llm = ChatOpenAI(model="gpt-4")

def llm_node(state: State) -> dict:
    response = llm.invoke(state["messages"])
    return {"messages": [response]}
\end{lstlisting}
\end{codebox}

\vspace{0.1cm}

{\footnotesize Works with \texttt{add\_messages} reducer: appends the response to \texttt{messages}.}
\end{frame}

\begin{frame}[fragile,t]{Node Patterns: Tool Execution}
\vspace{0.05cm}

\begin{codebox}[Tool Node]
\begin{lstlisting}
def tool_executor(state: State) -> dict:
    tool_name = state["tool"]
    tool_input = state["tool_input"]
    result = execute_tool(tool_name, tool_input)
    return {"tool_result": result}
\end{lstlisting}
\end{codebox}

\vspace{0.1cm}

{\footnotesize Bridges LLM requests and actual execution.}
\end{frame}

\begin{frame}[fragile,t]{Node Patterns: Routing Logic}
\vspace{0.05cm}

\begin{codebox}[Router Node]
\begin{lstlisting}
def router(state: State) -> dict:
    if state["score"] > 0.5:
        decision = "approve"
    else:
        decision = "reject"
    return {"decision": decision}
\end{lstlisting}
\end{codebox}

\vspace{0.1cm}

{\footnotesize Simple conditional logic. (For graph-level routing, use conditional edges-next section.)}
\end{frame}

\section{Edges}
\begin{frame}[t]{Edges: Connecting Nodes}
\setlength{\itemsep}{1pt}
\vspace{0.1cm}

\begin{conceptbox}[Three Types of Edges]
\setlength{\itemsep}{1pt}
\begin{itemize}
  \item \textbf{Normal edge:} Node A -> Node B (always)
  \item \textbf{Conditional edge:} Node A -> (B or C or ...) based on state
  \item \textbf{START / END:} Special entry/exit nodes
\end{itemize}
\end{conceptbox}

\vspace{0.15cm}

Edges define the \textit{control flow}. They answer: ``After this node, where do we go?''
\end{frame}

\begin{frame}[fragile,t]{Normal Edges: add\_edge()}
\vspace{0.1cm}

\begin{codebox}[Adding Edges]
\begin{lstlisting}
graph.add_edge("node_a", "node_b")
graph.add_edge(START, "first_node")
graph.add_edge("last_node", END)
\end{lstlisting}
\end{codebox}

\vspace{0.15cm}

\begin{tipbox}[START and END]
\texttt{START} = entry point. \texttt{END} = exit point. Graph begins at START, terminates at END.
\end{tipbox}

\vspace{0.15cm}

{\footnotesize A simple linear graph: START -> A → B → C → END}
\end{frame}

\section{Conditional Edges}
\begin{frame}[fragile,t]{Conditional Edges: add\_conditional\_edges()}
\vspace{0.05cm}

\begin{conceptbox}[Router Functions]
A router function reads state and returns the \textit{name} of the next node.
\end{conceptbox}

\vspace{0.1cm}

\begin{codebox}[Example]
\begin{lstlisting}
def route_decision(state: State) -> str:
    if state["decision"] == "approve":
        return "approve_node"
    else:
        return "reject_node"

graph.add_conditional_edges(
    "router",
    route_decision,
    {"approve_node": "approve_node",
     "reject_node": "reject_node"}
)
\end{lstlisting}
\end{codebox}

{\footnotesize Router returns node name string.}
\end{frame}

\begin{frame}[t,shrink=3]{START and END: Special Nodes}
\setlength{\itemsep}{0.5pt}
\vspace{0.02cm}

{\footnotesize
\begin{tabular}{p{2.5cm}p{7cm}}
\toprule
\textbf{Node} & \textbf{Purpose} \\
\midrule
\texttt{START} & Entry point. Initial state flows from here. \\
\texttt{END} & Exit point. Graph terminates, returns final state. \\
\bottomrule
\end{tabular}
}

\vspace{0.08cm}

\begin{tipbox}[Must Have]
\setlength{\itemsep}{0.5pt}
{\footnotesize
\begin{itemize}
  \item Edge from START
  \item Edge to END
\end{itemize}
}
\end{tipbox}
\end{frame}

\section{State Flow}
\begin{frame}[t]{How State Flows Through Nodes and Edges}
\vspace{0.1cm}

\begin{tikzpicture}[node distance=1.5cm, scale=0.95, every node/.style={font=\scriptsize}]
  \node[draw, fill=lcblue!30, rounded rectangle] (s1) {START\n(initial)};
  \node[draw, fill=lcgreen!30, rounded rectangle, right=of s1] (n1) {Node A};
  \node[draw, fill=lcgreen!30, rounded rectangle, right=of n1] (n2) {Node B};
  \node[draw, fill=lcred!30, rounded rectangle, right=of n2] (end) {END\n(return)};

  \path[->, thick] (s1) -- node[above, font=\tiny] {state} (n1);
  \path[->, thick] (n1) -- node[above, font=\tiny] {merged} (n2);
  \path[->, thick] (n2) -- node[above, font=\tiny] {final} (end);
\end{tikzpicture}

\vspace{0.15cm}

\begin{conceptbox}[Merging at Each Step]
\begin{enumerate}
\setlength{\itemsep}{1pt}
  \item Node receives \textit{full} state
  \item Node returns \textit{partial} dict
  \item Graph merges: \texttt{state.update(return\_dict)}
  \item Next node gets updated state
\end{enumerate}
\end{conceptbox}
\end{frame}

\section{Graph Lifecycle}
\begin{frame}[fragile,t,shrink=5]{The Graph Lifecycle}
\vspace{0.05cm}

\begin{codebox}[Step by Step]
\begin{lstlisting}
# 1. Define state
class State(TypedDict):
    messages: Annotated[list, add_messages]

# 2. Create graph
graph = StateGraph(State)

# 3. Add nodes
graph.add_node("llm", llm_node)

# 4. Add edges
graph.add_edge(START, "llm")
graph.add_edge("llm", END)

# 5. Compile
compiled = graph.compile()

# 6. Invoke
result = compiled.invoke({"messages": []})
\end{lstlisting}
\end{codebox}
\end{frame}

\section{Common Mistakes}
\begin{frame}[fragile,t,shrink=2]{Common Mistakes}
\vspace{0.05cm}

\begin{warningbox}[Mistake 1: Returning Full State]
\begin{lstlisting}
# WRONG
def node(state: State) -> State:
    return {"a": 1, "b": 2, "c": 3}  # All fields

# RIGHT
def node(state: State) -> dict:
    return {"a": 1}  # Only what changes
\end{lstlisting}
\end{warningbox}

\vspace{0.1cm}

\begin{warningbox}[Mistake 2: Forgetting START/END]
\begin{lstlisting}
# WRONG: no entry/exit
graph.add_edge("A", "B")

# RIGHT
graph.add_edge(START, "A")
graph.add_edge("B", END)
\end{lstlisting}
\end{warningbox}
\end{frame}

\begin{frame}[fragile,t]{Common Mistakes: Reducer Issues}
\vspace{0.1cm}

\begin{warningbox}[Mistake 3: Missing Reducer on List]
\begin{lstlisting}
# WRONG: messages will be replaced, not appended
class State(TypedDict):
    messages: list

# RIGHT: use add_messages reducer
class State(TypedDict):
    messages: Annotated[list, add_messages]
\end{lstlisting}
\end{warningbox}

\vspace{0.15cm}

{\footnotesize Without a reducer, returning \texttt{\{"messages": [new]\}} replaces the entire list!}
\end{frame}

\section{Summary}
\begin{frame}[t,shrink=2]{Summary: State, Nodes, Edges}
\setlength{\itemsep}{0.5pt}
\vspace{0.05cm}

\begin{conceptbox}[The Three Pillars]
\setlength{\itemsep}{1pt}
\begin{itemize}
  \item \textbf{State:} TypedDict schema + Annotated reducers
  \item \textbf{Nodes:} Functions (state -> dict updates)
  \item \textbf{Edges:} Routes (normal, conditional, START, END)
\end{itemize}
\end{conceptbox}

\vspace{0.15cm}

\begin{tipbox}[Key Patterns]
\setlength{\itemsep}{1pt}
\begin{itemize}
  \item \texttt{Annotated[list, operator.add]} for appending
  \item \texttt{add\_messages} for chat history
  \item Return partial dicts from nodes
  \item Define state before writing nodes
\end{itemize}
\end{tipbox}
\end{frame}

\begin{frame}[t,shrink=5]{Self-Check Questions}
\setlength{\itemsep}{0pt}
\vspace{0.05cm}

{\footnotesize
\begin{enumerate}
\setlength{\itemsep}{0.5pt}
  \item What is the purpose of a reducer in state?
  \item How does \texttt{operator.add} differ from default overwrite?
  \item What should a node function return?
  \item Why is \texttt{add\_messages} special for chat?
  \item What are the three edge types?
  \item What happens if you forget the START edge?
  \item How is state merged after a node executes?
\end{enumerate}

\vspace{0.05cm}

\textbf{Quick Answers:} (1) Merges updates (not replace). (2) Appends instead of replaces. (3) Partial dict of updates. (4) Smart chat merge logic. (5) Normal, conditional, START/END. (6) Graph won't start. (7) \texttt{state.update(node\_return)}.
}
\end{frame}

\end{document}
